% TOP_PROBLEM: {"id": 3374, "title": "Are all Mersenne Numbers with prime exponent square-free?", "category_id": 1, "category": {"id": 1, "name": "number_theory", "display_name": "Number Theory", "description": "Properties of integers, prime numbers, Diophantine equations.", "slug": "number-theory", "order_index": 1, "created_at": "2026-07-31T15:26:25.670Z"}, "status": "open", "problem_number": "OPG-59976", "set_id": 12}
% FIRST_POSTED: 2026-10-01
% ENTRY_KIND: statement_only
% UnsolvedMath ID 3374; source catalogue code OPG-59976
% !TeX program = lualatex
% Definition and sources only; this file is not an LLM solution attempt.
\documentclass[11pt,a4paper]{article}
\usepackage[margin=25mm]{geometry}
\usepackage{fontspec}
\setmainfont{lmroman10-regular.otf}[BoldFont=lmroman10-bold.otf,
 ItalicFont=lmroman10-italic.otf,BoldItalicFont=lmroman10-bolditalic.otf]
\usepackage{amsmath,amssymb,mathtools}
\usepackage{unicode-math}
\setmathfont{latinmodern-math.otf}
\AtBeginDocument{\let\setminus\smallsetminus}
\usepackage{xurl,hyperref}
\hypersetup{colorlinks=true,urlcolor=blue}
\setcounter{secnumdepth}{0}
\setlength{\parindent}{0pt}
\setlength{\parskip}{5pt}
\setlength{\emergencystretch}{3em}
\begin{document}
\section*{Are all Mersenne Numbers with prime exponent square-free?}\label{3374}
{\small UnsolvedMath ID 3374 · OPG-59976 · Number Theory · OpenGarden}
\subsection{Definitions and mathematical statement}
For a prime exponent $p$, set $M_p=2^p-1$. The question is whether
every such integer is square-free:
\[
 \forall\text{ primes }p,q,\qquad q^2\nmid 2^p-1.
\]
The exponent $p$ and a possible divisor $q$ are distinct roles and
need not be equal. Square-free means that every prime factor occurs
with multiplicity exactly one. It allows $M_p$ to be composite.
Restricting the exponent to primes is part of the problem, rather
than a conclusion to be established.

For $p=2$, one obtains $M_p=3$. For odd prime $p$, any prime divisor
$q$ of $M_p$ is odd and the multiplicative order of $2$ modulo $q$
is exactly $p$: the order divides $p$ and cannot be one. Thus
$p\mid q-1$. This necessary divisibility condition explains why the
two prime parameters cannot be treated as arbitrary divisors of a
general Mersenne number. A counterexample requires the stronger
congruence $2^p\equiv1\pmod{q^2}$ for an actual prime $p$.
Checking square-freeness for any bounded set of exponents leaves the
universal assertion beyond that bound untouched.
\subsection{Short English statement}
Must every Mersenne number with prime exponent have all its prime factors occurring only once?
\subsection{Sources}
\begin{itemize}
\item[S1] ULAM.AI and the UnsolvedMath Contributors, supplied problems.json, record 3374 (OPG-59976), OpenGarden collection. Catalogue identity and source transcription; CC BY 4.0.
\url{https://huggingface.co/datasets/ulamai/UnsolvedMath}.
\item[S2] Open Problem Garden, Are all Mersenne Numbers with prime exponent square-free?. Original problem and discussion; the mathematical formulation below is expanded from the supplied source record.
\url{http://www.openproblemgarden.org/op/are_all_mersenne_numbers_with_prime_exponent_square_free}.
\end{itemize}
\end{document}
